\chapter{Locally compact abelian groups}

Throughout this chapter, the letter G denotes, unless otherwise stated, a locally compact abelian group equipped with a Haar measure generally denoted by dx; for \(p \in [1, +\infty]\), the space \(L^{p}(G, dx)\) will simply be denoted \(L^{p}(G)\), and its norm will be denoted by \(f \mapsto \|f\|_{p}\). We identify \(L^{1}(G)\) with a subspace of \(\mathcal{M}^{1}(G)\) by the map \(f \mapsto f \cdot dx\). Recall that the support of the Haar measure is equal to G (INT, VII, §1, no. 1, remark 3); in particular (INT, III, §2, no. 2, proposition 9), the canonical mapping from the space \(\mathcal{K}(G)\) into \(L^{p}(G)\) is injective for \(p \in [1, +\infty]\). For \(p \neq +\infty\), we identify \(L^{p}(G)\) with a subspace of the space \(\widetilde{\mathcal{F}}(G; \mathbf{C})\), the space of classes of complex-valued functions on G defined and finite almost everywhere (INT, IV, §3, nos. 5, 6). In particular, the notation \(L^{1}(G) \cap L^{2}(G)\) denotes the intersection of \(L^{1}(G)\) and \(L^{2}(G)\) in this space. We denote by \(f \mapsto \widetilde{f}\) the involution on the involutive algebra \(L^{1}(G)\) (example 4 of I, p.~99); we have \(\widetilde{f}(x) = \overline{f(x^{-1})}\) for all \(x\) in G.

Recall (INT, V, §5, no. 3, th. 1) that if μ is a complex measure on a locally compact topological space X and if f is a locally μ-integrable function on X, then the measure ν with density f with respect to μ is denoted f·μ or simply fμ. A function g from X to C is essentially integrable for the measure ν if and only if gf is essentially integrable for \(\mu\); we then have

\[
(f \cdot \mu) (g) = \int_ {\mathrm{X}} g d (f \cdot \mu) = \int_ {\mathrm{X}} g f d \mu = \mu (g f).
\]

If G is a compact topological group, the normalized Haar measure on G is the unique Haar measure \(\mu\) on G such that \(\mu(\mathrm{G}) = 1\).

If \(X\) is a discrete topological space, the counting measure on \(X\) is the discrete measure \(\mu\) on \(X\) such that \(\mu(\{x\}) = 1\) for all \(x \in X\). If \(G\) is a discrete topological group, the counting measure on \(G\) is a Haar measure on \(G\).

We will also make use of the following two lemmas.

Lemma 1. --- Let X be a locally compact topological space and \(\mu\) a positive measure on X. Let \(x\) be an element of the support of \(\mu\). Let U be an open neighborhood of \(x\). There exists a function \(f \in \mathcal{K}_{+}(\mathrm{X})\) whose support is contained in U such that \(\int fd\mu = 1\).

By definition of the support of a measure (INT, III, §2, no. 2, def. 1), there exists a function \(g \in \mathcal{K}(\mathrm{X})\) supported in U such that \(\mu(g) \neq 0\) . We then have \(\mu(|g|) > 0\) since \(\mu\) is positive, and the function \(f = \mu(|g|)^{-1}|g|\) has the desired properties.

Lemma 2. --- Let G and H be topological groups with identity elements \(e_{\mathrm{G}}\) and \(e_{\mathrm{H}}\). Suppose the topological group G is Hausdorff. Let f be a morphism of topological groups from G to H. Suppose that for every neighborhood U of \(e_{\mathrm{G}}\) in G, there exists a neighborhood W of \(e_{\mathrm{H}}\) in H such that \(f^{-1}(\mathrm{W}) \subset \mathrm{U}\). Then the morphism f is injective and strict (TG, III, p.~16, def. 1).

The hypotheses imply that \(\operatorname{Ker}(f)\) is contained in every neighborhood of \(e_{G}\) in G, hence that \(\operatorname{Ker}(f) = \{e_{\mathrm{G}}\}\) since G is separated. The homomorphism from G to \(f(\mathrm{G})\) induced by f by passing to subspaces is then bijective; denote by \(g: f(\mathrm{G}) \to \mathrm{G}\) the inverse homomorphism. The hypotheses then imply that g is continuous at \(e_{H}\), therefore continuous (TG, III, p.~15, prop. 23).

